\documentclass[11pt,a4paper]{article}
\usepackage[margin=24mm]{geometry}
\usepackage{amsmath,amssymb,amsthm}
\setlength{\emergencystretch}{2em}
\newtheorem{theorem}{Theorem}
\newcommand{\gen}[1]{\langle#1\rangle}
\title{Two moving orbits in a singleton prefix}
\author{Human Compression project: L044 continuation}
\date{}
\begin{document}\maketitle
Retain the critical singleton setting with $n=2r\ge14$,
actual natural $G_0=\operatorname{Alt}(\Omega\setminus\{a\})$,
full ambient IP and the acquired strict prefix-growth proof.
For $4\le g<r$ write
\[
 P=G_{<g},\quad Y=\operatorname{supp}(P),\quad
 C=a^P,\quad H=G_{\{2,\ldots,g-1\}},\quad b=a^{\rho_0}.
\]
Assume $P$ is intransitive on $Y$ and $|Y|\le2g+1$.
In particular $Y$ is proper in $\Omega$, because
$|Y|\le2r-1<n$. The points $a,b$ are fixed by $H$.

\begin{theorem}
The group $H$ cannot have precisely two nontrivial orbits and
precisely the two fixed points $a,b$ within $Y$.
This exclusion does not require a paired-edge hypothesis.
\end{theorem}
\begin{proof}
Suppose the stated orbit decomposition holds, and name the
nontrivial orbits $U,V$. Strict prefix growth gives
$|C|\ge g+1\ge5$. Thus $C$ meets at least one of $U,V$,
and contains every $H$-orbit it meets. Since $a,b\in C$,
meeting both would make $P$ transitive on $Y$.
Relabel so that
\[
                  C=\{a,b\}\mathbin{\dot\cup}V,
                       \qquad E=U.                       \tag{1}
\]
These are the two $P$-orbits in $Y$, so $\rho_1$ preserves
both sets.

The mark $\rho_0$ centralizes $H$, and therefore permutes
its nontrivial orbits equivariantly. If it exchanges $U,V$,
then it is fixed-point-free on their union; if it preserves
one of them, its restriction is either the identity or a free
involution. This last assertion follows because a centralizing
permutation with a fixed point fixes a whole transitive orbit.

If $\rho_0$ exchanges $U,V$, or preserves both and is free
on both, then $\operatorname{supp}(\rho_0)=Y$.
The mark $\rho_1$ preserves $Y$, and every mark of index at
least two centralizes $\rho_0$ and preserves its support.
Thus $Y$ is a proper nonempty $G$-invariant set, a contradiction.
If $\rho_0$ is the identity on both $U,V$, it is just
$(a\ b)$ and is odd, again impossible. Consequently it is
free on exactly one of the two orbits and fixes the other
pointwise.

If it is free on $V$ and fixes $E$, then
$\operatorname{supp}(\rho_0)=C$ by (1).
The mark $\rho_1$ preserves $C$, and every later mark
centralizes $\rho_0$. Thus $C$ is a proper $G$-invariant set,
a contradiction. The sole remaining case is
\[
 \rho_0|_E\text{ is free},\qquad \rho_0|_V=1.             \tag{2}
\]

Put $P_1=\gen{\rho_0,\rho_1}$. Strict growth makes
$x=b^{\rho_1}\ne b$. Since $\rho_1$ fixes $a$ and
preserves the two sets in (1), $x\in V$.
The $P_1$-orbit of $a$ is exactly $\{a,b,x\}$:
on it the marks are $(a\ b)$ and $(b\ x)$.
Actual natural $G_0$ and full IP give
\[
       (P_1)_a=P_1\cap G_0=\gen{\rho_1}.
\]
Orbit--stabilizer therefore gives $|P_1|=3\cdot2=6$.
In particular $P_1\cong S_3$ and
$(\rho_0\rho_1)^3=1$ in the actual group, including
its action outside $\{a,b,x\}$.

The set $V\setminus\{x\}$ is preserved by both marks:
$\rho_0$ fixes it, and $\rho_1$ preserves $C$ and has
already paired $b$ with $x$ and fixed $a$.
On this set $\rho_0\rho_1=\rho_1$ has order dividing
two. Its order also divides three by the displayed relation,
so both marks fix $V\setminus\{x\}$ pointwise.
Every point outside $Y$ is fixed as well.

Every $P_1$-orbit in $E$ has size dividing six. It cannot
have size one or three, because $\rho_0$ is free there and
an involution on an odd set has a fixed point. Hence these
orbits have size two or six. We have proved that $E$ is
exactly the union of all $P_1$-orbits of sizes two and six
in the whole ambient set; the sole other nontrivial orbit is
$\{a,b,x\}$ of size three.
Every mark of index at least three centralizes $P_1$, so
preserves this union $E$. The first two marks preserve it
by (1), and $\rho_2\in H$ preserves it as an $H$-orbit.
Thus every mark preserves $E$, the final contradiction.
\end{proof}

\paragraph{Consequences and limits.}
In the numerical budget for a small singleton prefix, this theorem
eliminates the entire row $f=2,c=2$, including both possible
support sizes $2g$ and $2g+1$. The remaining rows
$f=3,c=2$ and $f=2,c=3$, both with support $2g+1$, require
separate arguments. The theorem above assumes its literal
orbit decomposition directly; it does not depend on the status
of the separate budget packet.

\paragraph{Attribution.}
The actual singleton point stabilizer, strict prefix growth and
critical range are acquired. The support-invariance split and the
use of the actual order-six first prefix to distinguish all exterior
orbits are this continuation. No faithfulness of a general orbit
projection is used. The source context is
Cameron--Fernandes--Leemans--Mixer, \emph{Highest rank of a polytope
for $A_n$}, Propositions~5.28--5.32.
No finite census, formal declaration, whole-original recovery or
question credit is added.
\end{document}
